%! Intercepts, Zeros, and Extrema — Unit 2, Lesson 2.4 — Student Packet Cover Sheet
\documentclass[10pt]{article}
\usepackage{atda-article}
\usepackage{atda-boxes}
\usepackage{ltablex}
\keepXColumns

\begin{document}

% Full-bleed royal banner
\begin{tikzpicture}[remember picture, overlay]
  \fill[royal] (current page.north west)
    rectangle ([yshift=-0.9in]current page.north east);
\end{tikzpicture}
\vspace{-0.8in}

\begin{center}
  {\color{white}\textbf{\LARGE Algebra, Trigonometry, and Data Analysis}}\\[4pt]
  {\color{mistmid}\large\textbf{Unit 2: Functions, Transformations, and Graph Analysis}}\\[6pt]
  {\color{white}\textbf{\large Lesson 2.4 \quad Intercepts, Zeros, and Extrema}}
\end{center}

\vspace{0.22in}
\namedateperiod
\vspace{0.18in}

\begin{learningtargetbox}
\small
\begin{enumerate}[leftmargin=1.5em, itemsep=5pt]
  \item \ldots{} find the \textbf{$y$-intercept} and the \textbf{$x$-intercepts} of a graph, and give
        the \textbf{zeros} of the function --- saying each time whether my answer is an input, an
        output, or a point.
  \item \ldots{} name the intervals where a graph is \textbf{increasing}, \textbf{decreasing}, or
        \textbf{constant}, and write them in interval notation using the \emph{input} axis.
  \item \ldots{} identify the \textbf{absolute maximum} and \textbf{absolute minimum} of a function
        over its domain, reporting a \emph{value} and a \emph{location} every time.
  \item \ldots{} explain why a function may have no zeros, no absolute maximum, or two places where
        the same extreme value happens.
  \item \ldots{} interpret every one of these characteristics in the situation the graph models.
\end{enumerate}
\end{learningtargetbox}

\vspace{0.14in}

\begin{tocbox}
\small
\renewcommand{\arraystretch}{1.5}
\begin{tabularx}{\linewidth}{@{} c l X r @{}}
\rowcolor{mist}
\textbf{\#} & \textbf{Component} & \textbf{Description} & \textbf{Score} \\
\midrule
1 & Warm-Up        & Two ways of asking ``where is the output zero?'' & \blank{1.2cm} \\
2 & Guided Notes   & Intercepts and zeros, rising and falling, and the highest and lowest points of
                     a graph                                        & \blank{1.2cm} \\
3 & Group Activity & Reading a gym's busiest hour, a start-up's break-even months, and two claims
                     that use the wrong axis                        & \blank{1.2cm} \\
4 & Exit Ticket    & One creek, five characteristics, and a classmate who answered in the wrong
                     units                                          & \blank{1.2cm} \\
5 & Homework       & The bike-share station, a cooling cup of coffee, and a look ahead
                                                                    & \blank{1.2cm} \\
\midrule
  & \multicolumn{2}{r}{\textbf{Total}} & \blank{1.2cm} \\
\end{tabularx}
\end{tocbox}

\vspace{0.14in}

\begin{remindbox}
\small
This lesson covers \texttt{AFDA.AF.2b} (intervals where a function is increasing, decreasing, or
constant), \texttt{AFDA.AF.2c} (the location \emph{and} value of the absolute maximum and minimum
over the domain), and \texttt{AFDA.AF.2d} (zeros and intercepts). Everything here is read
\textbf{off a graph} --- that is not a shortcut, it is what this course asks of you. One sentence runs
the whole packet: \textbf{before you write a number down, say which axis it came from.} An interval
of \emph{inputs} and a stretch of \emph{outputs} can be the same two numbers and mean completely
different things.
\end{remindbox}

\end{document}
